\documentclass{amsart}
\usepackage{graphicx} % Required for inserting images

\title{Thesis log}
\author{Riley Wood}
\date{\today}


\usepackage[style = ieee, backend = biber]{biblatex}
\addbibresource{../bibliography.bib}

\begin{document}

\maketitle
\section{Week $2$}

\subsection{What I learned this week}
Reviewed Kripke structures, abstract interpretation, our favorite Galois connection, and widening operators (in the general sense).
Overall not a great week for new content, I spent a lot of time on reminding myself what is even going on with abstract interpretation.

\subsection{What I wrote this week}
A bit of an introduction to formal verification. 
I don't think it's very good, but it should give me a framework to build upon down the line.
Maybe struggling with precision when writing this, since I want it to be really accessible.
It might be better in the future to start with definitions and formalism, then build narrative around that.

\subsection{Questions}
\begin{itemize}
  \item Is my notion of program state correct? (no one seems to bother formalizing this anywhere)
  \item What is going on with the zonotope definition that $\cite{Zonotopes}$ gives?
  \item what is the two variables per inequality domain?
  \item What is going on with the approximations in Greg's paper? (Riemann approximations. Can we do better for polytopes?)
  \item What is there to be said about Rice's theorem?
  \item Kleene algebras show up when talking about methods for building state diagrams. What are these?
\end{itemize}

\subsection{Next week plan}
I think that this week I may have been too focused on writing narrative and studying fundamentals. 
Going forward maybe focus on more on formalism, and build narrative later.

Revisiting ToPAZ pretty closely might be worthwhile, but I have spent a fair bit of time with it already.

I will spend a lot of time on $\cite{PreciseWidening}$. 
34 pages on how to do polyhedra widenings! Yes please. 
This article seems like a good central point to do some research. 
It cites some things I know will be important (Chernikova!) and gets cited fairly frequently (recently by a paper on selective widening and two variables per inequality)

\section{Week $1$}

\subsection{What I learned this week}
\begin{itemize}
  \item Read \cite{DomainTheory}, a paper looking to use algebraic effects and domain theory to give a calculus and abstract machine for functional logic programming. 
They base this off of \textbf{Call By Push Value} (CBPV). 
It is a little more domain theoretic than I maybe am planning to get with this project, but it does fit into a pattern of papers from this year's popl working on algebraic effects and CBPV.
  \item 
Read\cite{Hyperfunctions}, a paper which looks to make a generalization for continuations and algebraic effects in the concurrent setting. 
They build a framework for understanding this pattern that shows up in a bunch of concurrent functional programming contexts (and also with church numerals(!). 
There could be a project looking to find of implement this in the wild
  \item 
Read \cite{Hasegawa-Thielecke}, another paper interested in CBPV, this time looking to do a bunch of category theory to include both Call by Name and Call by Value interpretation into one category or something like that.
The issue is that these operations break associativity, so they use Duploids (what are these?) to solve their problems.
  \item 
Read \cite{SemanticsToSyntax}, a pretty type-theory focused paper which tries to reflect some structure from comprehension categories back into the type systems they support.
They offer some prohects to extend their work, tryign to show their syntax is complete (?) and trying to compare their work to bicategorical type theory.
  \item 
Read \cite{WhatIsMonoid}, a paper which looks to build monoidal structures that can describe everything we see in CS that looks like a monoid but doesnt work out.
In particular they are interested in generalizing relative monads (whatever those are) and CBPV . 
Once again we get a discussion of associativity being broken by CBN and CBV.
They solve this with bi-skew multicategories.
It seems like their is a project to compare all these different ways of grappling with how CBPV breaks things. Maybe some research could emerge from that.
This paper is also pretty generous with future work. 
Syntactic treatment of their monoid, trying to work on a theorem for representing bi-skew multicategories, and exploring what monoid actions look like (maybe in a multi\emph{act}agory - whatever that means).

\end{itemize}

\subsection{What I wrote this week}
Brief summaries of the papers referenced above, some ideas for projects that we can discuss in our next meeting.

\subsection{Questions}
I tried very hard not to dive into every definition and calculation in these papers, so I have questions about pretty much every programming language and category theoretic term referenced above.

\subsection{Next week plan}
Before our next meeting (M-TH), continue lit review, focusing slightly more on abstract interpretation review.

In our meeting: decide between available project options (will discuss them in as much detail as I am able then) .

After our meeting: look for "tentpole" papers and allow a bit more time to dive into common definitions (and WRITING the common definitions).

\printbibliography[]

\end{document}
